% =========================================================
% Set Algebras
% =========================================================

\subsection{Set Algebras}

\begin{exposition}
Closure rules become mathematical structures when they are packaged together.
A set algebra is the finite Boolean closure package for families of subsets.
It is the first structure in this restart where the family is required to
behave like a small universe of sets.
\end{exposition}

\begin{definitionbox}{Definition (Set Algebra)}
\begin{definition}[Set Algebra]\label{def:set-algebra}
Let \(X\) be a set.  A \emph{set algebra} on \(X\) is a family
\(\mathcal A\subseteq\mathcal P(X)\) such that:
\begin{enumerate}[label=(\roman*)]
\item \(X\in\mathcal A\);
\item if \(A\in\mathcal A\), then \(X\setminus A\in\mathcal A\);
\item if \(A_1,\dots,A_n\in\mathcal A\), then
      \(\bigcup_{k=1}^{n}A_k\in\mathcal A\), with the empty finite union
      interpreted as \(\varnothing\).
\end{enumerate}
\end{definition}
\end{definitionbox}

\begin{remark*}[Standard quantified statement]
\[
\mathcal A\text{ is a set algebra on }X
\iff
X\in\mathcal A,\quad
\mathcal A\text{ is complement-closed},\quad
\mathcal A\text{ is closed under finite unions}.
\]
\end{remark*}

\begin{remark*}[Interpretation]
A set algebra is a family that admits the ambient set, complements, and finite
unions.  The previous closure-duality theorem explains why finite
intersections do not have to be listed separately in this definition.
\end{remark*}

\begin{remark*}[Exposition]
Many books call this a \emph{field of sets}.  In this chapter, \emph{set
algebra} is the preferred term because it points directly at the finite
Boolean operations.
\end{remark*}

\LeanFormalizes{def:set-algebra}{lra-lean}{LRA.VolumeI.Set.Families}{SetAlgebra}{pending}

\begin{dependencies}
\begin{itemize}
  \item \hyperref[def:closed-under-set-operation-vocabulary]{Closed Under Set Operations}
  \item \hyperref[thm:pairwise-closure-implies-finite-closure]{Pairwise Closure Implies Finite Closure}
  \item \hyperref[thm:complement-union-intersection-closure-duality]{Complement-Union-Intersection Closure Duality}
  \item \hyperref[def:power-set]{Power Set}
  \item \hyperref[def:complement]{Complement}
  \item \hyperref[def:union]{Union}
\end{itemize}
\end{dependencies}

\begin{theorembox}{Theorem (Set Algebra Equivalent Definitions)}
\begin{theorem}[Set Algebra Equivalent Definitions]\label{thm:set-algebra-equivalent-definitions}
Let \(\mathcal A\subseteq\mathcal P(X)\).  The following packages are
equivalent:
\begin{enumerate}[label=(\roman*)]
\item \(\mathcal A\) is a set algebra: \(X\in\mathcal A\), \(\mathcal A\) is
      closed under complements, and \(\mathcal A\) is closed under finite
      unions.
\item \(\varnothing\in\mathcal A\), \(\mathcal A\) is closed under
      complements, and \(\mathcal A\) is closed under finite unions.
\item \(X\in\mathcal A\), \(\mathcal A\) is closed under complements, and
      \(\mathcal A\) is closed under finite intersections.
\item \(\mathcal A\) contains \(X\), is closed under complements, is closed
      under finite unions, and is closed under finite intersections.
\end{enumerate}
\hyperref[prf:set-algebra-equivalent-definitions]{Go to proof.}
\end{theorem}
\end{theorembox}

\begin{remark*}[Standard quantified statement]
\[
\begin{aligned}
&\mathcal A\text{ is a set algebra on }X\\
&\iff
\varnothing\in\mathcal A,\ \mathcal A\text{ is complement-closed, and }
\mathcal A\text{ is closed under finite unions}\\
&\iff
X\in\mathcal A,\ \mathcal A\text{ is complement-closed, and }
\mathcal A\text{ is closed under finite intersections}.
\end{aligned}
\]
\end{remark*}

\begin{remark*}[Interpretation]
These equivalent definitions say that a set algebra can be introduced with
several axiom packages.  Once complements are available, finite union closure
and finite intersection closure determine each other.
\end{remark*}

\LeanFormalizes{thm:set-algebra-equivalent-definitions}{lra-lean}{LRA.VolumeI.Set.Families}{SetAlgebraEmptyAxiomPackage}{pending}
\LeanFormalizes{thm:set-algebra-equivalent-definitions}{lra-lean}{LRA.VolumeI.Set.Families}{SetAlgebraFiniteIntersectionPackage}{pending}
\LeanFormalizes{thm:set-algebra-equivalent-definitions}{lra-lean}{LRA.VolumeI.Set.Families}{SetAlgebraFiniteBooleanClosurePackage}{pending}
\LeanFormalizes{thm:set-algebra-equivalent-definitions}{lra-lean}{LRA.VolumeI.Set.Families}{SetAlgebraEquivalentDefinitions}{pending}
\LeanFormalizes{thm:set-algebra-equivalent-definitions}{lra-lean}{LRA.VolumeI.Set.Families}{SetAlgebraIffEmptyAxiomPackage}{pending}
\LeanFormalizes{thm:set-algebra-equivalent-definitions}{lra-lean}{LRA.VolumeI.Set.Families}{SetAlgebraIffFiniteIntersectionPackage}{pending}
\LeanFormalizes{thm:set-algebra-equivalent-definitions}{lra-lean}{LRA.VolumeI.Set.Families}{SetAlgebraIffFiniteBooleanClosurePackage}{pending}

\begin{dependencies}
\begin{itemize}
  \item \hyperref[def:set-algebra]{Set Algebra}
  \item \hyperref[def:empty-set]{Empty Set}
  \item \hyperref[thm:complement-laws]{Complement Laws}
  \item \hyperref[thm:complement-union-intersection-closure-duality]{Complement-Union-Intersection Closure Duality}
\end{itemize}
\end{dependencies}

\begin{theorembox}{Theorem (Set Algebra Closure Consequences)}
\begin{theorem}[Set Algebra Closure Consequences]\label{thm:set-algebra-closure-consequences}
Let \(\mathcal A\) be a set algebra on \(X\).  Then:
\begin{enumerate}[label=(\roman*)]
\item \(\varnothing\in\mathcal A\);
\item \(\mathcal A\) is closed under complements;
\item \(\mathcal A\) is closed under finite unions;
\item \(\mathcal A\) is closed under finite intersections;
\item if \(A,B\in\mathcal A\), then \(A\setminus B\in\mathcal A\);
\item if \(A,B\in\mathcal A\), then \(A\triangle B\in\mathcal A\).
\end{enumerate}
\hyperref[prf:set-algebra-closure-consequences]{Go to proof.}
\end{theorem}
\end{theorembox}

\begin{remark*}[Standard quantified statement]
\[
\mathcal A\text{ is a set algebra on }X
\to
\varnothing\in\mathcal A,\quad
\mathcal A\text{ is complement-closed},\quad
\mathcal A\text{ is closed under finite unions, finite intersections, differences, and symmetric differences}.
\]
\end{remark*}

\begin{remark*}[Interpretation]
After this theorem, a set algebra can be used as a finite Boolean workspace:
all of the ordinary finite set operations stay inside the algebra.
\end{remark*}

\LeanFormalizes{thm:set-algebra-closure-consequences}{lra-lean}{LRA.VolumeI.Set.Families}{SetAlgebraContainsEmpty}{pending}
\LeanFormalizes{thm:set-algebra-closure-consequences}{lra-lean}{LRA.VolumeI.Set.Families}{SetAlgebraClosedUnderComplements}{pending}
\LeanFormalizes{thm:set-algebra-closure-consequences}{lra-lean}{LRA.VolumeI.Set.Families}{SetAlgebraClosedUnderFiniteUnions}{pending}
\LeanFormalizes{thm:set-algebra-closure-consequences}{lra-lean}{LRA.VolumeI.Set.Families}{SetAlgebraClosedUnderFiniteIntersections}{pending}
\LeanFormalizes{thm:set-algebra-closure-consequences}{lra-lean}{LRA.VolumeI.Set.Families}{SetAlgebraClosedUnderPairwiseDifferences}{pending}
\LeanFormalizes{thm:set-algebra-closure-consequences}{lra-lean}{LRA.VolumeI.Set.Families}{SetAlgebraClosedUnderPairwiseSymmetricDifferences}{pending}

\begin{dependencies}
\begin{itemize}
  \item \hyperref[def:set-algebra]{Set Algebra}
  \item \hyperref[thm:set-algebra-equivalent-definitions]{Set Algebra Equivalent Definitions}
  \item \hyperref[thm:difference-laws]{Difference Laws}
  \item \hyperref[thm:symmetric-difference-laws]{Symmetric Difference Laws}
  \item \hyperref[thm:complement-union-intersection-closure-duality]{Complement-Union-Intersection Closure Duality}
\end{itemize}
\end{dependencies}
